\cotefile{83}
\scanpage{79}
% loc centre, diagram rows 1–4
\diagram{\begin{tikzcd}[ampersand replacement=\&,execute at end picture={
\begin{scope}[shift={(current bounding box.center)},x=.014cm,y=.014cm]
\node at (0,200) {\texttt{Totalité} $(\mathrm{IV}'_{1})$};
\node[anchor=east] at (-105,68) {\texttt{simplicité} $(\mathrm{IV}'_{2})$};
\node[anchor=west] at (103,68) {\texttt{Ordre} $(\mathrm{IV}'_{6})$};
\node[anchor=east] at (-105,-65) {\texttt{Unité} $(\mathrm{IV}'_{3})$};
\node[anchor=west] at (103,-65) {\texttt{Causalité} $(\mathrm{IV}'_{5})$};
\node at (0,-200) {\texttt{Structure}\,\unsure{.}$(\mathrm{IV}'_{4})$};
\draw (-64,170)--(-141,94) (53,175)--(153,87)
(-160,45)--(-161,-42) (149,49)--(148,-45)
(-141,-108)--(-52,-172) (37,-172)--(138,-91);
\draw (-102,77)--(92,75) (-94,-66)--(93,-67)
(-20,151)--(-137,-31) (17,160)--(135,-31)
(-120,36)--(-15,-144) (23,-141)--(134,39);
\draw (0,49)--(0,-44) (-47,31)--(47,-31) (-48,-33)--(47,35);
\draw[thin] (-50,-39)--(44,22);
\draw[line width=.9pt] (-141,297) .. controls (-119,168) and (-76,-93) .. (-19,-207) .. controls (-12,-239) and (-2,-267) .. (11,-279);
\node[inner sep=0pt] at (-25,-194) {$\deleted{\vphantom{I}}$};
\end{scope}}]
\phantom{\rule{11.4cm}{8.4cm}}
\end{tikzcd}}
% archive sheet number: 78

\scanpage{81}
% loc 1, top heading and date
\uline{Azumaya rang} 4 \mspage{I} Avril \unsure{85|81|86}.

% loc 2–4
1. Soient $k$ anneau, $A$ algèbre d'Azumaya sur $k$ de rang 4 (donc loc. iso. à $M_2(k)$). On a alors
% loc 5–6
\[
\left\{\begin{aligned}
\operatorname{Tr}: A&\longrightarrow k & k\text{-linéaire}\\
\det: A&\longrightarrow k & \text{multiplicatif, quadratique}
\end{aligned}\right. \tag*{(1)}
\]
% loc 7
tels que $\forall u\in A$, on ait
% loc 8
\[ u^2-\operatorname{Tr}(u).u+\det(u)=1 \tag*{(2)} \]
% loc 9
(Hamilton-Cayley). De plus, on a
% loc 10
\[ \det(u+v)=\det u+\det v+(\operatorname{Tr}u\operatorname{Tr}v-\operatorname{Tr}(uv)) \tag*{(3)} \]
% loc 11
(i.e.
% loc 12
\[ (u,v)\longmapsto\unsure{\psi|\Psi|\Phi}(u,v)\overset{\text{dfn}}{=}\operatorname{Tr}u\operatorname{Tr}v-\operatorname{Tr}(uv) \]
% loc 13
est la forme bil. \unsure{ass.} à la forme quad. $\det$).

% loc 14–19, remark
\uline{\unsure{Remarque|Remarques}} On peut s'intéresser aux $k$-Alg. quelconques $A$ pour lesquelles on a (1) (2) (3). Si $A$ est proj. de t.f. disons et 2 inv. dans $k$, alors $(1)(2)\Longrightarrow(3)$, car il suffit de le voir pour $u=v$ (deux fonctions bil. \uline{sym} qui coïncident pour $u=v$ coïncident)\unsure{.|,}
% loc 19–22, cancelled continuation and restart
\deleted{et alors cela se \inserted{résulte de} réduit à une vérification dans les $k[u]=k\oplus k.u$ quadratiques sur $k$ \unsure{...})} La relation pertinente est
% loc 23
\[ (\operatorname{Tr}u)^2-\operatorname{Tr}(u^2)=2\det u, \tag*{(4)} \]
% loc 24
et elle résulte de (2) en \unsure{s'écrisant|l'écrivant} sous forme
% loc 25, between bottom lines
\[ u^2=\operatorname{Tr}(u)u-\det(u) \tag*{(2')} \]
% loc 26, bottom edge
et portant cette valeur de $u^2$ dans (4), compte tenu
% archive sheet number: 80

\scanpage{82}
% loc 1
de la première des relations
% loc 2–3
\[ \left\{\begin{aligned}\operatorname{Tr}1&=2\\ \det 1&=1\end{aligned}\right. \tag*{(5)} \]
% loc 4–6
\deleted{les quelles}\inserted{Dont la 2°} se déduit \deleted{\unsure{est}} de $(2')$ en y faisant $u=1$ (et $\det 1=1$ est donné comme connu).

% loc 7–8
En polarisant (2), on trouve pour $u,v\in A$ (tenant compte de 3)
% loc 9
\[ uv+vu-\operatorname{Tr}(\unsure{v|u})u-\operatorname{Tr}(u)\unsure{v|u}+(\operatorname{Tr}u\operatorname{Tr}v-\operatorname{Tr}(uv))1 \tag*{(6)} \]
% loc 10
soit
% loc 11
\[ uv+vu=\operatorname{Tr}v)u+(\operatorname{Tr}u)v+(\operatorname{Tr}(uv)-\operatorname{Tr}u\operatorname{Tr}v)1 \tag*{$(6')$} \]
% loc 11–20, oblique note from left edge
% slanted
\margin{\uline{NB} Cette relation $(6')$ (pour la seule donnée de la forme $\operatorname{Tr}$, sans dét ni \unsure{Ham. Cayley|Ham.-Cayley}.) équivaut \unsure{``au} fait que $u\mapsto\bar u'=\operatorname{Tr}(u).1-u$ est un ``antiautomorphisme'' (\unsure{en vertu de|en vertu} \deleted{de carré l'\unsure{id.} identité}). Ceci dit, pour les \unsure{deux} ``formes'' $\operatorname{Tr}$, \unsure{l'une|la} quadratique $u\mapsto\det u$, les deux relations de $(2)\ (3)\ (6)$ ($\unsure{\Longleftrightarrow|\Rightarrow}6''$) impliquent la troisième.}
% loc 12–13
\deleted{\uline{Corollaire} On en déduit aussitôt} Cela \deleted{signif} implique que mod
% loc 14
\[ k.1+ku+kv, \]
% loc 15
$uv$ et $vu$ sont \deleted{\unsure{opposés}} opposés
% loc 16
\[ vu\equiv uv\qquad\bmod\quad k1+ku+kv. \]
% loc 17–20, proposition
\uline{Proposition} Soit $A$ avec (1) (2) (3) (d'où $(6')$) Alors $\forall u,v\in A$, \deleted{$k+ku+kv+k.\unsure{v|uv}$} le sous $k$-module engendré par $1,u,v,uv$ est une sous-Algèbre.

% loc 21, bottom
Il suffit de voir que \deleted{le produit de}
% archive sheet number: 81

\scanpage{83}
% loc 1–4, struck block
\deleted{deux \deleted{éléments} parmi les éléments $1,u,v,uv$ est dans ce sous-module. Il suffit de \deleted{le} voir pour deux tels éléments distincts de 1. \deleted{\unsure{Il|On}} faut regarder}
% loc 5–7, struck multiplication table
\[\begin{aligned}
\deleted{u.u}&\deleted{=\det(u)1+(\operatorname{Tr}u)u+\unsure{\det}}\\
\deleted{u.v}&\deleted{=0.1+0.u+0.v+1uv}\\
\deleted{u.uv}&\deleted{=}
\end{aligned}\]
% loc 8, crossed list
\[ \deleted{\underline{u^2},\underline{v^2},uvuv,\underline{uv},u^2v,\underline{vu},vuv,uvu,uv^2} \]
% loc 9–12, struck calculations
\deleted{Or \unsure{$u^2$}}
\[\begin{aligned}
\deleted{u^2}&\deleted{=\operatorname{Tr}(u).u-\det(u).1}&&\deleted{\text{((2) pour }u}\\
\deleted{v^2}&\deleted{=\operatorname{Tr}(v).v.-\det(v)1}&&\deleted{\text{(2) pour }v}\\
\deleted{uv}&\deleted{=1.uv}&&\deleted{\text{tautol}}\\
\deleted{vu}&\deleted{=-1.uv+\operatorname{Tr}(v)u+\operatorname{Tr}u\,v+(\operatorname{Tr}(uv)-\operatorname{Tr}u\operatorname{Tr}v)1}
\end{aligned}\]
\deleted{(relation $(6')$)}
% loc 13–18, new proof
Ce sous-module $V$ est stable par multiplication \inserted{\uline{à gauche}} par \unsure{$1$}, $u,v$, \deleted{$uv$ \unsure{.} \illegible{3}} \inserted{(donc aussi par} $uv$, sans compter 1). Donc il suffit de voir que $ux,vx\in V$, si $x\in V$. Il suffit de prendre pour $x$ un des \unsure{gén|gèn} $1,u,v,uv$. Le cas 1 \unsure{étant|est} trivial, il reste
% loc 19
\[ \underline{u^2},\underline{uv},\underline{u^2}v,\;vu,\underline{v^2},vuv=\deleted{\unsure{v(uv)}}(vu)v \]
% loc 20–23
Pour $u^2,v^2$ on utilise H.C. (3), pour $uv$ c'est tautologique, pour $u^2v$ cela résulte de l'expression de $u^2$ H.C., pour $vu$ on utilise $(6')$, qu'on écrit
% loc 24, overwritten formula
\[ \deleted{vu=-\unsure{uv}+\operatorname{Tr}(v).u+{}} \]
% loc 25
\[ vu=(\operatorname{Tr}(uv)-\operatorname{Tr}u,\operatorname{Tr}v)1+\operatorname{Tr}(v)u+\operatorname{Tr}(u)v-uv. \tag*{$(6'')$} \]
% loc 26–29
Enfin pour \deleted{$(vu)v$} $vuv=(vu)v$, on multiplie la relation $(6'')$ à droite par $v$, en utilisant la relation H.C. pour $v$ (pour se débarrasser des termes en $v^2$).
% archive sheet number: 82

\scanpage{84}
% loc 1–5
Finalement, on trouve une ``table de multiplication'' pour une $k$-algèbre libre de rang 4 (dont $V$ sera quotient) \unsure{explicite|explicitée} en termes de 5 éléments de $k$
% loc 6–8, display and notes at right
\[
\left\{\begin{aligned}
 b&=\operatorname{Tr}u,\ c=\det u\deleted{\unsure{,a}}&
 &\text{(permet de décrire } k[u]\inserted{=k\oplus k.u}\\
 &&&\quad \inserted{\text{(si }k\cap k.u=\{0\}\text{)}}\text{)}\\
 b'&=\operatorname{Tr}v,\ c'=\det v\deleted{\unsure{=}}&&\text{(id pour } k[v]\text{)}\\
 t&=\operatorname{Tr}(uv)
\end{aligned}\right. \tag*{(7)}
\]
% loc 9–12
avec comme base $\underline{1},\underline{u},\underline{v},\underline{uv}$ (les soulignés indiquent qu'on considère $1,u\ \unsure{..}$ comme éléments de la nouvelle algèbre). On appelle cette Algèbre
% loc 13
\[ M(\deleted{\unsure{a}}b,\deleted{\unsure{b}}c;\deleted{\unsure{a'}}b',\deleted{\unsure{b'}}c';t) \tag*{(8)} \]
% loc 14–16
On vérifie qu'elle existe (indépendamment de la donnée de $A$) pour tout choix de \deleted{$b,c,$}
% loc 17
\[ (b,c,b',c',t)\in k^5 \tag*{(\unsure{4|9})} \]
% loc 18–19
En fait, on a une alg. canonique \unsure{de} rang 4 sur
% loc 20
\[ \mathbb Z[B,C,B',C',T] \]
% loc 21
où $B,C,B',C',T$ sont des indéterminées.

% loc 22–24
\uline{Vérification} On écrit une table de multiplication explicite, en utilisant H C pour \unsure{$u,v$|$u$ et $v$} et (6) alias $(6')$ (pour $(u,v)$),
% archive sheet number: 83

\scanpage{85}
% loc 1–5
et on vérifie l'associativité. De plus, on prolonge $\operatorname{Tr}$ à $V$ par linéarité, et $\det$, \deleted{par} d'abord sur $\underbrace{k[\underline u]}_{A_1}\subset M$ et $\underbrace{k[\underline v]}_{A_2}\subset M$ (où c'est \deleted{$\det(u'$} la norme des algèbres quadratiques, $\operatorname{Tr}$ y étant la Trace) puis sur
% loc 6–7
\[ M=k[u]\oplus k[v]\qquad\text{(}\inserted{\text{iso}}\ \text{du }\text{pt de vue}\ \text{additif)} \]
% loc 8
par
% loc 9
\[ \det(u_1^{\deleted{\unsure{2}}}+v_1)=\underbrace{\det u_1}_{N_{A_1/k}(u_1)}+\underbrace{\det v_1}_{N_{A_2/k}(v_1)}+\operatorname{Tr}u_1\operatorname{Tr}v_1-\operatorname{Tr}(u_1v_1) \tag*{(10)} \]
% loc 10–13
Ceci dit, ``on \deleted{voit}\inserted{vérifie} que'' $\det$ est quadratique et multiplicatif, que l'on a la relation de Hamilton-Cayley et la relation (3).

% loc 14–18
En fait, je n'ai pas eu à faire la vérification, car \deleted{\unsure{je sai|je sais}} il suffit de la faire (principe de prolongement des id alg) \deleted{pour} en supposant
% loc 19
\[ \delta\overset{\text{dfn}}{=}b^2-4c\ \text{inv.},\quad \delta'\overset{\text{dfn}}{=}b'^2-4c'\ \text{inv.} \]
% loc 20–22
(i.e. $A_1,A_2$ étales sur $k$), et on \deleted{sait} voit que, \deleted{pa}\inserted{pour $b,c,b',c'$ donnés ainsi,} l'on \deleted{peut} réaliser \deleted{(pour alors (pour pour) le \unsure{système|septime|répartition} $(b,c$}
% archive sheet number: 84

\scanpage{86}
% loc 1–7
le système $(b,c,b',c',t)$ ``pour presque tout \deleted{choi} $t$'' \inserted{et loc fpqc} par deux éléments \inserted{$u,v$} \unsure{de dans|dedans|donnant} l'alg de Matrices \inserted{$A=M_2(k)$}. En effet, \deleted{partant} \deleted{de \unsure{une|$u,v$}}\inserted{quitte} à se localiser pour que $A_1,A_2$ soient \unsure{scindies|scindées} $(\simeq k\times k)$, on peut les réaliser comme plongées dans $A$, i.e. par $u\in A$, $v\in A$.
% loc 8–10
Bien sûr, on ne \deleted{pas de \unsure{en|ce}} peut préjuger de la valeur de $\operatorname{Tr}uv$. Mais \deleted{considérons, pour}\inserted{on a}

% loc 11–14, lemma
\uline{Lemme} Soient $u,v\in A\ (\simeq M_2(k))$ qui \deleted{soient} engendrent $k[u],k[v]$ étales rang 2 sur base $\operatorname{Spec}k$. Considérons pour $\unsure{\theta|\Theta}$ variable dans $A^*$
% loc 15
\[ \operatorname{Tr}(u\theta(v)) \qquad \theta(v)\overset{\text{dfn}}{=}\theta v\theta^{-1} \]
% loc 16
donc on trouve morphisme de schémas
% loc 17
\[ \deleted{\unsure{M_2|M}}\ \check{\mathbb V}(A)\longrightarrow\underbrace{\mathbb E^1_k}_{\text{espace affine}=\operatorname{Spec}k[T]} \]
% loc 17–18, oblique note to left
% slanted
\margin{\uline{NB} $\theta\unsure{\longrightarrow|\longmapsto}\theta(v)$ est épi sur les conjugués à $v$,}
% loc 18
Celui-ci est \uline{dominant} --- 

% loc 19–23
En fait, on trouve une fonction rationnelle en $b,c,b',c',t$, ayant comme dénominateur $t^2-4cc'$ $\bigl(=\deleted{\det}\operatorname{Tr}uv^2-4\det(uv)=\delta(uv)$, discriminant de $uv\bigr)$
% archive sheet number: 85

\scanpage{87}
% loc 1–2
(l'invariant \uline{anharmonique} de $k[u]$ et $k[v]$ dans alg. d'Azumaya)
% loc 3
\[ \deleted{\unsure{\text{de}}}\quad \rho=f(b,c,b',c';t)\deleted{\unsure{a}}=\frac{P(b,c,b',c',t)}{Q(\text{---})} \]
% loc 4–7
telle que la condition néc. et suffisante pour que \deleted{$M$}$(b,c,b',c',t)$ \deleted{soit alg. d'Azumaya}\inserted{se réalise loc. fpqc par alg. d'Azumaya}, \deleted{c'est que} (en supposant déjà $\delta\delta'$ inversible), c'est que $\rho\inserted{\text{partout}}\ne2$,
% loc 8
\[ \text{i.e.}\qquad P-2Q\inserted{\text{partout}}\ne2 \]
% loc 9–18
En fait, quand ils sont réalisés dans alg. d'Azumaya, \deleted{ils définissent} alors $A_1,A_2$ définissent \deleted{deux div} chacun un diviseur de degré 2\deleted{, et} sur \deleted{\unsure{l'id}} le schéma projectif associé à $A$, et si leurs supports sont disjoints, il y a un inv. anharmonique \deleted{``\unsure{défini|fois}''}$\lambda$ (qui n'est défini que localement \inserted{sur $\operatorname{Spec}k$} et mod involution $\lambda\mapsto\lambda^{-1}$) et l'invariant (défini sur $k$)
% loc 19
\[ \bigl(\qquad \rho=\lambda+\deleted{\frac{1}{\lambda}}\lambda^{-1}\qquad \text{d'où}\quad \lambda^2-\rho\lambda+1=0 \]
% loc 20–21
(et comme $\lambda$ partout $\ne0,1$, on doit avoir $\rho$ partout $\ne2\ \ldots$
% archive sheet number: 86

\scanpage{88}
% loc 1–5
J'ai fait le calcul de $\rho$ (en termes de $b,c,b',c',t$) dans le cas où $b=\operatorname{Tr}u$ et $b'=\operatorname{Tr}v$ sont nuls, \deleted{on trouve}\inserted{et} supposant que $u,v$ plongés dans A Alg d'Az. On trouve alors
% loc 6–7
\[
\left.\begin{aligned}
\rho&=2\frac{t^2+4cc'}{t^2-4cc'}\\
\text{d'où}\quad \frac{\rho}{2}-1&=8\frac{cc'}{t^2-4cc'}
\end{aligned}\right|\ \text{si } b=b'=0
\]
% loc 6–7, right-hand NB
\margin{\uline{NB} $t^{\unsure{1|2}}-4cc'=\delta(uv)$ discriminant de $uv$}
% loc 8–10
On va en déduire le calcul dans le cas général. [Il suffit de le faire sur $\mathbb Z[\frac12][B,C,B',C',T]$.) . On trouve
% loc 11–12, display (11) and annotations beneath equalities
\[
\rho\underset{\substack{\downarrow\\\text{OK si 2 inv.}\\\text{dans }k}}{=}
2\frac{(2t-bb')^2+\delta(u)\delta(v)}{(2t-b\deleted{\unsure{b}}b')^2-\delta(u)\delta(v)}
\underset{\substack{|\\\text{OK}\\\text{ds ts les}\\\text{cas}}}{=}
\frac{2t^2-2bb't+\bigl(b^2b'^2-2(b^2c'+b'^2c)+8cc'\bigr)}
{t^2-bb't+(b^2c'+b'^2c-4cc')} \tag*{(11)}
\]
% loc 13, bracketed specialization
[ En \uline{car. 2}, cela donne $\displaystyle\rho\underset{\text{car. 2}}{=}\frac{(bb')^2}{t^2-bb't+(b^2c'+b'^2c)}$. ]
% loc 14–15
\[
\rho_0\overset{\unsure{\uparrow}}{=}\rho-2
=\frac{b^2b'^2-4(b^2c'+b'^2c)+16cc'}{t^2-bb't+(b^2c'+b'^2c-4cc')}
=\frac{A(b,b',c,c')}{B(t,b,b',c,c')\deleted{\unsure{t}}}
\tag*{\deleted{(\unsure{12|**})}}
\]
\margin{invariant anh. rectifié}
% loc 16–17
Donc
\[ \rho B(t,b,b',c,c')\deleted{\unsure{\rho}}-A(b,b',c,c')=0 \]
% archive sheet number: 87

\scanpage{89}
% loc 1–3
$A_1$ et $A_2$ \inserted{diagonales} engendrent $M_2$. Cela se réalise au moins loc. fpqc (en passant à la \deleted{solution} rev. des solutions de
% loc 4
\[ \lambda^2-\rho\lambda+1=0\quad\deleted{\unsure{\rho}} \]
% loc 5–8
d'autre part, à cause de $b^2-4c$ \unsure{inv.|\deleted{inv.}}, on peut trouver loc. fpqc (passant au rev. des solutions de $x^2-bx+c=0$), $u$ dans $A_1$ correspondant à
% loc 9
\[ \operatorname{Tr}u=b,\quad \det u=c \]
% loc 10
et item $v$ dans $A_2$
% loc 11
\[ Tv=b',\quad \det v=c'. \]
% loc 12
Soit \deleted{$t'$}
\[ \tau=\operatorname{Tr}uv \]
% loc 13–15
alors $\tau$ satisfait la relation analogue à $(11')$ avec $t$ remplacé par $\tau$. Or en tant qu'équation en $t$
% loc 16
\[ t^2-bb't+\left(b^2c'+b'^2c-4cc'-\frac{\delta\delta'}{\rho_0}\right) \]
% loc 17–19
(où $\delta=b^2-4c,\ \delta'=b'^2-4c'$), il y a \deleted{tout juste} (sur un corps) \deleted{au plus} 2 solutions, \deleted{soient} $t,t'$ (\unsure{évt.|év!} égales) satisfaisant
% loc 20
\[ t+t'=bb' \]
% loc 21–22, bottom
Si $t=\tau$, on a gagné. Sinon \inserted{$(\tau=t')$} on remplace $u$ par
% archive sheet number: 88

\scanpage{90}
% loc 1–4, upper display and gloss
\[
\left\{\begin{aligned}
\rho_0=\rho-2
&=\frac{\delta\deleted{(u)}\delta'\deleted{(v)}}{t^2-bb't+(\delta c'+\delta'c+4\deleted{\unsure{b}}cc')}\\
&=\frac{(b^2-4c)(b'^2-4c')}{\underbrace{t^2-bb't+b^2c'+b'^2c-4cc'}_{Q=Q(t;b,b',c,c')}}
\end{aligned}\right.\qquad
\begin{aligned}\text{où }\delta&=\delta(u)=b^2-4c\\ \delta'&=\delta(v)=b'^2-4c'\end{aligned}
\tag*{$(11')$}
\]
% loc 5
\unsure{$1$}

On trouve donc :

% loc 6–7, proposition
\uline{Proposition} Soient $k$ un anneau, et $b,c,b',c',t\in k$, tels que
% loc 8
\[ (b^2-4c)(b'^2-4c')(t^2-bb't+b^2c'+b'^2c-4cc') \]
% loc 9–15
soit inv. dans $k$. Alors $M[b,c,b',c',t]$ est une alg. d'Azumaya. (Et $k[\underline u],k[\underline v]$ en sont des sous-algèbres quadratiques étales, \inserted{``en position générale'' i.e \unsure{engendrent|engendrant} \unsure{la dite|de cette} Alg.}) Dans tous les cas, c'est une algèbre associative, \deleted{et} satisfaisant les relations (2) (3) pour $\operatorname{Tr},\det$ définis comme plus haut\ldots.

% loc 16–20, lower-left marginal OBS
% slanted
\margin{OPS $k$ corps alg. clos, car c'est une question \unsure{géom.} fibre par fibre.}
% loc 16–21, proof
\uline{Dém} On regarde $\rho_0$ comme défini plus haut $(11')$. Comme $\rho_0$ \deleted{\unsure{étant}} inv i.e $\rho$ partout $\ne2$, cela correspond à une position relative de tores dans un $\operatorname{SL}(2)$ (ou de sous alg. quadratiques \inserted{$A_1,A_2$} étales dans un $M_2$) en \uline{position générale}. -- Cela signifie aussi que
% archive sheet number: 89

\scanpage{91}
% loc 1–2
\[
u'=(\operatorname{Tr}u)1-u=b.1-u
\qquad\left\{\begin{aligned}\operatorname{Tr}u'&=\operatorname{Tr}u\\ \det u'&=\deleted{\unsure{\operatorname{Tr}v}}\det u\end{aligned}\right.
\]
% loc 3–5
alors
\[\begin{aligned}
\tau'&=\operatorname{Tr}u'v=\operatorname{Tr}\deleted{\unsure{u}}(b-u)v=b\operatorname{Tr}v)-\operatorname{Tr}(uv)\\
&=bb'-\tau\deleted{\unsure{(uv)}}
\end{aligned}\]
donc \deleted{\unsure{$t'$}} $\tau'=t$, On gagne !

% loc 6
Que se passe-t-il si
% loc 7
\[ (b^2-4c)(b'^2-4c')\in k^{\unsure{\times|*}} \tag*{$(12_1)$} \]
% loc 8
\deleted{mais} \inserted{Si} $M(b,c,b',c',t)$ est azumayen, je dis que
% loc 9
\[ \underbrace{t^2-bb't+b^2c'+b'^2c-4cc'}_{Q(t;b,c,b',c')\deleted{\unsure{t}}}\in k^*\quad ? \tag*{$(12_2)$} \]
% loc 10–11
\deleted{Je dis qu'alors $M(b,c,b',c',t)$ \uline{n'est pas} \unsure{Azumayen|Azumaya|Azumayenne}.} En effet, \deleted{si c'était}
% loc 12–17
Comme $A_1=k[\underline u]$ et $A_2=k[\underline v]$ \deleted{\unsure{q}} sont étales, et en position générale, il \deleted{les formules} y aurait par la théorie générale un invariant $\rho_0$, relié à $b,c,b',c',t$ par $(11')$, récrit sous forme
% loc 18
\[ \deleted{\rho_0(t^2-bb't}\qquad \rho_0Q=\delta\delta' \]
% loc 19
Comme $\delta\delta'\in k^*$, on doit avoir $Q\in k^*$

% loc 20–21
\uline{Remarque} Le discriminant de $Q$ (en tant que polynôme en $t$) est
% archive sheet number: 90

\scanpage{92}
% loc 1
\[ \Delta(Q)=(bb')^2-4(b^2c'+b'^2c-4cc') \]
% loc 2
On trouve
% loc 3
\[ \Delta(Q)=(b^2-4c)(b'^2-4c')=\delta\delta' \tag*{(13)} \]
% loc 4–7
donc en fait, si \deleted{\unsure{(}} (comme on a supposé) $\delta\delta'\in k^{\unsure{\times|*}}$, alors $\unsure{Q(k)|Q(b)|\Delta(Q)}\in k^{\times}$ i.e ``racines distinctes''. Cela signifie aussi que \deleted{$t\ne bb'-t$ i.e $2t\ne bb'$}
% loc 8–10, cancelled lines
\[ \deleted{2\operatorname{Tr}(uv)=} \]
\[ \deleted{\delta\delta'\in k^{\unsure{1}*}\Longrightarrow 2t\ne bb'\ \text{partout}}\tag*{\deleted{(14)}} \]
\deleted{i.e $2t-bb'\in k^\times$}
% loc 11–14
\deleted{\unsure{Bref}} Le rev. des solutions de
\[ Q(t;b,c,b',c')\deleted{\unsure{t}}=0 \]
(en $t$, pour $b,c,b',c'$ fixés) est étale (quadratique) sur $S=\operatorname{Spec}k$.

% loc 15–19, left margin
% slanted
\margin{\deleted{\uline{A vérifier}. \unsure{Immédiat} sur corps, par ce qui précède}}
% loc 15–20, corollary
\uline{Cor.} (?) Soient $u,v$ dans alg. d'Azumaya \inserted{$A$}, tel que $\delta(u)\delta(v)\in k^*$, \deleted{\unsure{si|$k$}} et $u,v$ partout \inserted{lin. indép. mod 1.} \inserted{et que} $Q(u,v)=Q(t;b,c,b',c')=0$, \deleted{alors} il \deleted{$u,v$ appart. à une \unsure{même}} tels que $u,v$ appartiennent à une \deleted{\unsure{même}} sous-Alg. de Borel (qui n'est autre alors que $k\oplus ku\oplus kv$) $\Longleftrightarrow uv\in k.1+ku+kv$
% archive sheet number: 91

\scanpage{93}
% loc 1–5, multiplication table
\diagram{\scriptsize\setlength{\arraycolsep}{2pt}$\begin{array}{c|c|c|c|c}
 &1&\underline u&\underline v&\underline{uv}\\ \hline
1&\underline 1&\underline u&\underline v&\underline{uv}\\
\underline u&\underline u&-c+b\underline u&\underline{uv}&-c\underline v+b\underline{uv}\\
\underline v&\underline v&(t-bb')\underline1+b'\underline u+b\underline v-\underline{uv}&-c'\underline1+b'\underline v&-bc'\underline1+c'\underline u+t\underline v\\
\underline{uv}&\underline{uv}&-b'c\underline1+t\underline u+c\underline v&-c'\underline u+b'\underline{uv}&t\underline{uv}-cc'\\ \hline
\end{array}$}

% loc 6–7
En prenant les traces, on a la matrice
% loc 8–11, cancelled matrix
\[
\left(\begin{array}{c|c|cc}
\deleted{2}&\deleted{\unsure{b}\ b}&\deleted{b'}&\deleted{t}\\
\deleted{b}&\deleted{-2c+b^2}&\deleted{t}&\deleted{-b'\unsure{c}+bt}\\
\deleted{b'}&\deleted{\unsure{2(bb')}\,t}&\deleted{-2c'+b'^2}&\deleted{-\unsure{2}bc'+\unsure{bc'}+b't}\\
\deleted{t}&\deleted{\underline{-\unsure{2}b'c+bt}+\unsure{c}}&\deleted{-bc'+b't}&\deleted{t^2-cc'}
\end{array}\right)
\]
% loc 12–15, second matrix
\[
\begin{pmatrix}
2&b&b'&t\\
b&\deleted{\unsure{4}}b^2-2c&t&bt-b'c\\
b'&t&b'^2-2c'&b't-bc'\\
t&bt-b'c&b't-bc'&t^2-cc'
\end{pmatrix}
\]
% loc 12–15, note at left
\margin{dét. de $\alpha,\beta\mapsto\operatorname{Tr}(\alpha\beta)$ p.r. à la base \unsure{partielle|particulière} $1,u,v,uv$)}
% loc 16–18, lower matrix and caption
$\cdot$ dét de $\alpha,\beta\longmapsto\operatorname{Tr}(\alpha)\unsure{\operatorname{Tt}|\operatorname{Tr}}(\beta)$
\[
\begin{pmatrix}
4&2b&2b'\\2b&b^2&bb'\\\unsure{2b|2b'}&bb'&b'^2
\end{pmatrix}
\]
% archive sheet number: 92

\scanpage{94}
% loc 1, page identifier
PU \mspage{\handwritten{\unsure{65}}}

% loc 2–6, upper crossed paragraph
\deleted{Cette abstraction est \deleted{apparente en}\inserted{\handwritten{inhérente à}} tout mot \deleted{sans exception} du langage\inserted{ sans exception}\inserted{\handwritten{\unsure{,|.}}} \deleted{\unsure{et}} \handwritten{C'}est \deleted{même}\handwritten{bien} la \inserted{\handwritten{aussi une}} condition \deleted{sine qua non}\handwritten{nécessaire}, pour que le mot soit bien le patrimoine culturel d'un groupe ayant une continuité dans le temps, et non la propriété d'une personne\inserted{\handwritten{ ni d'un groupement de personnes}}. Il en est ainsi même pour les mots les plus élémentaires\inserted{\handwritten{ de tous}}, comme ``mère'', ``père'', ``manger'', ``boire'', ``soleil'' etc.}
% loc 4, left margin
\margin{\handwritten{\unsure{d très|très}}}
% loc 7–12
\deleted{Entre de tels mots-\deleted{\illegible{1}}\handwritten{base}\deleted{, \unsure{en}} et des mots comme ``famille'', voire ``groupe'', ``peuple'', ``nation'', ``gouvernement'', ``politique'', ``philosophie'', ``abstraction'', on sent \deleted{\unsure{une} une différence qualitative, \unsure{marquée}, une différence de ``degré d'abstraction''}\inserted{\handwritten{une ``distance''}}, comparable à celle qui existe entre les \deleted{mots ``élémentaires''}\handwritten{``mots-base''} eux-mêmes, et les ``objets concrets'' \deleted{(au sens le plus strict\unsure{,} de tantôt)} auxquels ils s'appliquent.}

% loc 13–21, second paragraph
\deleted{Pour le dire autrement: le processus d'abstraction, qui \deleted{est}\inserted{semble être} (au point de vue de la formation de \uline{concepts}, porteurs de \uline{sens}) \deleted{l'opération essentielle}\handwritten{comme l'âme} dans la formation progressive d'un langage de plus en plus complexe, de plus en plus ``ramifié'', \deleted{et} apte à saisir \deleted{\unsure{les}}des replis et \deleted{\unsure{les}}des nuances \deleted{les} de plus en plus délicats\deleted{\unsure{e}} des choses de l'Univers -- ce processus est c\,u\,m\,u\,l\,a\,t\,i\,f . Une fois un de ses pas accomplis, le nouveau concept, s'incarnant dans un mot nouveau, finit par faire partie des choses familières\inserted{ et ``concrètes''}, au même titre, voire plus encore, que ces autres choses que nous avions qualifiées\inserted{\handwritten{ tantôt}} de \deleted{\unsure{``}}v\,r\,a\,i\,m\,e\,n\,t concrètes. \deleted{C'est sûr que pour un mathématicien}}
% loc 22–34, paragraph sign before Ainsi
\margin{\handwritten{\unsure{\P}}}
\deleted{Ainsi, le concept de n\,o\,m\,b\,r\,e (entier, disons), historiquement ou structuralement parlant, est\inserted{\handwritten{ visiblement}} d'un degré d'abstraction élevé (bien plus élevé qu'un mot-concept comme ``deux'' ou ``trois'' par exemple, \handwritten{le}qu\handwritten{el} représente\deleted{\unsure{nt}} \deleted{\unsure{une qualité|une qualités}} tout ce qu'il y a de terre-à-terre). Cela n'empêche qu'il n'y a pas un seul mathématicien, \inserted{sûrement,} à qui cette notion n'appara\handwritten{îsse} comme extrêmement \handwritten{``}concrète\handwritten{''} \deleted{par contre} -- bien plus concrète disons que ``feu'' ou ``terre'', \deleted{qui} désign\handwritten{an}t des choses avec lesquelles il a plus ou moins perdu contact (à supposer qu'il l'ait jamais eu). Au niveau d'abstraction où travaillent \deleted{la plupart}\inserted{beaucoup} des mathématiciens, le nombre \deleted{entier} apparaît aussi \inserted{\handwritten{concret (et en même temps, aussi}} lointain!\handwritten{)}, depuis les hauteurs où il \deleted{fait}\handwritten{trace} ses virevoltes, que \deleted{apparaîtra}\inserted{\handwritten{le}} serait le garde-manger de \handwritten{s}a maison familiale, à celui qui parlerait de l'économie française ou du marché mondial des céréales.}

% loc 35–40, final paragraph
\deleted{Tout ceci pour illustrer qu\deleted{e}\handwritten{'} \deleted{la notion d'``abstraction'',} à toutes fins pratiques et une fois réalisé que \uline{toute} pensée et \uline{tout} langage sont ``abstraits'' (par rapport aux ``objets'' au sens courant du terme), \inserted{\handwritten{la notion d'``abstraction'',}} est avant tout une \deleted{\illegible{1}} n\,o\,t\,i\,o\,n r\,e\,l\,a\,t\,i\,v\,e . Un discours, un langage, une pensée, une thé\unsure{orie|oire} sont ``plus abstraits'' ou ``plus concrets'' qu'un\inserted{\handwritten{ (ou une)}} autre. Quand l'un et l'autre\inserted{\handwritten{ (ou l'une et l'autre)}} appartiennent à des domaines pas trop éloignés de la}
% archive sheet number: 93

\scanpage{95}
% loc 1–7
Il faudrait prouver que ça implique qu'il y a un $x\in k.1\oplus ku\oplus kv$, partout non nul, qui est orth. (pour $\operatorname{Tr}(xy)$) à $1,u,v$ -- il faut donc calculer le \deleted{déterminant} discriminant de $\operatorname{Tr}$ par rapport à la base $1,u,v$, on trouve
% loc 8–11
\[
\begin{aligned}
&\text{discriminant}\\[-2pt]
&\text{de }\operatorname{Tr}\text{ pour}\\[-2pt]
&\text{la base }1,u,v\\[-2pt]
&\text{de }k1\oplus ku\oplus kv
\end{aligned}
=-2\,Q(t;b,c,b',c')
\]
% loc 12–14, bracketed addition
[ et le discriminant \uline{réduit} de la forme \uline{quadratique} $\det$ est
\[ -Q(t;b,c,b',c')\qquad\biggr] \]
% loc 12–15, left margin
% slanted
\margin{? Vérifier si c'est ça \uuline{Oui}}
% loc 15–17
Donc le fait que c'est nul \deleted{signifie}\inserted{implique} que \deleted{\unsure{les}} $\exists\ \unsure{\theta|\Theta}$ orthogonal à $1,u,v$, partout non nul et \uuline{isotrope}

% loc 18–19
Donc \deleted{que}
\[ \det\theta=\operatorname{Tr}\theta=0 \]
% loc 20–22
Or \deleted{l'une des x} les droites isotropes dans l'hyperplan $\operatorname{Tr}\theta=0$ -- correspondent aux Borels, \deleted{comme les} on gagne !
% archive sheet number: 94
